% &%%%%%%%%%%%%%%%%%%%%%%%%%%%%%%%%%%%%%%%%%%%%%%%%%%%%%%%%%%%
% &%%                                                       %%
% &%%                    Discussion                         %%
% &%%                                                       %%
% &%%%%%%%%%%%%%%%%%%%%%%%%%%%%%%%%%%%%%%%%%%%%%%%%%%%%%%%%%%%

% !===========================================================

\graphicspath{{5_Images/8_Discussion/}}                         % Chemin vers le stockage des images du chapitre
\makeatletter
\def\input@path{{5_Images/8_Discussion/}}                       % Pareil que au dessus
\makeatother

% !===========================================================

\chapter{Discussion}\label{chap:Discussion}                     % Ajout d'un nouveau chapitre

\begin{chapabstract}                                            % Environnement ajoutant un résumé de chapitre
    \lettrine[slope=0.5em,nindent=0.5em]{C}{eci}
    est un chapitre où vous pourrez présenter l’avancée de vos travaux. Vous y expliquerez les idées principales développées, les démarches ou méthodes utilisées, ainsi que les résultats et observations obtenus. L’objectif est de donner au lecteur une vue d’ensemble de ce qui est traité dans le chapitre, tout en mettant en évidence les points importants et les contributions à votre étude. Ce texte sert de guide pour orienter le lecteur avant de détailler les analyses, discussions et conclusions dans les sections suivantes.
\end{chapabstract}

\minitoc%                                                       % Ajout du mini-ToC
\newpage                                                        % Commence une nouvelle page pour le chapitre

% !===========================================================

% ?-----------------------------------------------------------
\section{Section 1}\label{sec:1_1}
% ?-----------------------------------------------------------

\lipsum[1]%

% ~-----------------------------------------------------------
\subsection{Sous-section 1}\label{subsec:1_1.1}
% ~-----------------------------------------------------------

\lipsum[1-2]%

% ------------------------------------------------------------
\subsubsection{Sous-sous-section 1}\label{subsubsec:1_1.1.1}
% ------------------------------------------------------------

\lipsum[1-4]%

% ------------------------------------------------------------
\subsubsection{Sous-sous-section 2}\label{subsubsec:1_1.1.2}
% ------------------------------------------------------------

\lipsum[1-5]%

% ------------------------------------------------------------
\subsubsection{Sous-sous-section 3}\label{subsubsec:1_1.1.3}
% ------------------------------------------------------------

\lipsum[1-2]%

% ~-----------------------------------------------------------
\subsection{Sous-section 2}\label{subsec:1_1.2}
% ~-----------------------------------------------------------

\lipsum[1-4]%

% ?-----------------------------------------------------------
\section{Section 2}\label{sec:1_2}
% ?-----------------------------------------------------------

\lipsum[1-2]%

% ~-----------------------------------------------------------
\subsection{Sous-section 1}\label{subsec:1_2.1}
% ~-----------------------------------------------------------

\lipsum[1-2]%

% ------------------------------------------------------------
\subsubsection{Sous-sous-section 1}\label{subsubsec:1_2.1.1}
% ------------------------------------------------------------

\lipsum[1-4]%

% ------------------------------------------------------------
\subsubsection{Sous-sous-section 1}\label{subsubsec:1_2.1.2}
% ------------------------------------------------------------

\lipsum[1-3]%

% ~-----------------------------------------------------------
\subsection{Sous-section 2}\label{subsec:1_2.2}
% ~-----------------------------------------------------------

\lipsum[1-3]%

% ------------------------------------------------------------
\subsubsection{Sous-sous-section 1}\label{subsubsec:1_2.2.1}
% ------------------------------------------------------------

\lipsum[1]%

% ------------------------------------------------------------
\subsubsection{Sous-sous-section 1}\label{subsubsec:1_2.2.2}
% ------------------------------------------------------------

\lipsum[1-2]%

% ~-----------------------------------------------------------
\subsection{Sous-section 2}\label{subsec:1_2.3}
% ~-----------------------------------------------------------

\lipsum[1-4]%

% ?-----------------------------------------------------------
\section{Section 3}\label{sec:1_3}
% ?-----------------------------------------------------------

\lipsum[1-5]%

% ?-----------------------------------------------------------
% ? BILAN
% ?-----------------------------------------------------------

\begin{chapbilan}
    Cette conclusion résume les points principaux abordés dans le chapitre et met en évidence les contributions ou enseignements obtenus. Elle permet de rappeler les résultats ou observations les plus importants, de souligner leur signification dans le cadre global de l’étude, et d’indiquer les perspectives ou questions ouvertes qui seront développées dans les chapitres suivants. L’objectif est de donner au lecteur une synthèse claire et structurée avant de passer à la suite du manuscrit.
\end{chapbilan}
